\documentclass[a4paper,11pt,fleqn]{article}%fleqn pour aligner les equations à gauche
\usepackage{/home/rweye/Nextcloud/Documents/Overleaf/Styles/Cours}


\newcommand{\chnum}{Chapitre 4}
\newcommand{\chtitle}{Le noyau de l'atome}
\newcommand{\dtype}{Cours}

\begin{document}
\maketitle

\section{Modèle de l'atome}
	\subsection{Constituants de l'atome}
	L'atome est constitué d'un noyau chargé positivement entouré d'un cortège d'électrons chargés négativement. Les composants du noyau sont appelés nucléons. Il existe deux sortes de nucléons : les neutrons et les protons.
	
	\subsection{Caractéristiques des constituants de l'atome}
	Le noyau est composé de :
	\begin{itemize}
		\item \textbf{protons} : chargés positivement de charge $+e = \SI{1,602E-19}{C}$ et ayant une masse $m_{proton}=\SI{1,673E-27}{kg}$
		\item \textbf{neutrons} : électriquement neutres et ayant une masse $m_{neutron}=\SI{1,675E-27}{kg}$
	\end{itemize}
	Les \textbf{électrons} sont situés autour du noyau. Ils ont une charge négative égale à $-e$ et une masse $m_{electron}=\SI{9,109E-31}{kg}$.\\
	Les masses des neutrons et des protons étant trop proches, on les considérera égales pour la suite des exercices. La masse moyenne d'un nucléon étant alors $m_{nucléon}=\SI{1,67E-27}{kg}$.\\
	La masse d'un électron est environ \num{2000} fois plus faible que celle d'un nucléon.
	
	\subsection{Notation symbolique du noyau d'un atome}
	
\end{document}
